\section{开放问题：泛化、黑箱与数据效率}

\subsection{泛化和记忆之间的张力}

Manning 在这一讲里反复强调一个直觉：大模型的成功并不意味着它们已经像人类一样理解世界。很多收益来自于海量训练数据和强大的统计模式记忆。

\begin{figure}[H]
\centering
\includegraphics[width=0.95\textwidth]{slides-images/slide-023.jpg}
\caption{从符号系统出发看 AI：符号、逻辑、编程语言和人类语言都属于“符号化的世界表征”\protect\footnotemark}
\end{figure}
\footnotetext{Slides 讨论了符号系统与认知科学、工程系统之间的关系。}

一个重要的对比是数据效率：
\begin{itemize}
    \item 人类往往可以从极少量示例中学会一个新技能；
    \item Transformer 在很多设置下仍然更依赖大量数据；
    \item 在某些低数据场景里，LSTM 甚至表现出更好的归纳偏置。
\end{itemize}

\begin{importantbox}{这里真正值得追问的不是“能不能拟合”}
更关键的问题是：
\begin{enumerate}
    \item 模型是否学到了可复用的结构？
    \item 它们是在泛化，还是在检索相似模式？
    \item 这些能力在什么条件下会失效？
\end{enumerate}
\end{importantbox}

\subsection{黑箱与可解释性}

神经网络仍然常常表现为黑箱。我们通常知道最终分数，却不知道模型内部到底学了什么。于是，解释性研究变得重要：

\begin{itemize}
    \item \textbf{Mechanistic interpretability}：从内部电路角度解释模型行为；
    \item \textbf{Causal abstraction}：寻找更高层次的可解释机制；
    \item \textbf{Neuron / feature analysis}：观察某些神经元或特征是否对应特定语言现象。
\end{itemize}

这类研究不是纯理论兴趣，而是因为我们真的需要知道模型是怎样成功、怎样失败的。

\subsection{本章小结}

大模型带来了能力跃迁，但也暴露出三个长期问题：数据效率不够、人类式泛化仍未实现、内部机制仍不透明。

%% ========== 语言学的角色 ==========

